\exo \textit{[calculer]}

 \begin{enumerate}
\item Calculer mentalement l'image de chaque nombre par la fonction racine carrée.

a)~9~ \hspace{0.8cm }b)~0~ \hspace{0.8cm}c)~49~~ \hspace{0.8cm}d)~$10^8$~~ \hspace{0.8cm}e)~$\dfrac{16}{25}$.

\item Calculer mentalement l'image de chaque nombre par la fonction carré.

a)~$-8$ \hspace{0,5cm} b)~$8$ \hspace{0,5cm} c)~$0,3$ \hspace{0,5cm} d)~$\dfrac{3}{4}$ \hspace{0,5cm} e)~$-\dfrac{4}{5}$ \hspace{0,5cm} f)~$10^3$ \hspace{0,5cm}

\item Calculer mentalement l'image de chaque nombre par la fonction cube.

a)~$3$ \hspace{0,5cm} b)~$-1$ \hspace{0,5cm} c)~$\dfrac{1}{2}$ \hspace{0,5cm} d)~$10^-2$ \hspace{0,5cm} 
\end{enumerate}

\exo \textit{[Raisonner, calculer]}

On considère la fonction racine carrée et $\mathcal{C}_f$ sa courbe représentative.

Parmi les points suivants, quels sont ceux qui appartiennent à $\mathcal{C}_f$? Justifier la réponse.

a)~$A(408,4441~;~20,21)$~ \hspace{0.8cm }b)~$B(6~;~5,1)$.

\medskip

\exo \textit{[Raisonner]}

En utilisant le sens de variation des fonctions de référence,  comparer les deux nombres donnés. 

\medskip

a)~$\sqrt{12}$~et~$\sqrt{10}$~\hspace{0.6cm}b)~${(-6,18)}^2$~et~$(-6,08)^2$~\hspace{0.6cm}c)~$(-4)^3$~et~$(-\pi)^3$~\hspace{0.6cm}d)~$(10^{-3})^2$~et~$(10^{-1})^2$ ~\hspace{0.6cm} e)~$(-4)^3$~et~$(-\pi)^3$.

\exo \textit{[Raisonner]}

On donne ci-dessous les courbes de trois fonctions de référence.

\begin{minipage}{0.45\linewidth}
Courbe 1:

\begin{center}
\psset{xunit=0.18cm,yunit=0.62cm,arrowsize=2pt 2}
\begin{pspicture}(-1,-0.4)(17.5,4.4)
  \psgrid[subgriddiv=5,gridlabels=0pt,gridcolor=gray!35,subgridcolor=gray!18](0,0)(17,4)
  \psaxes[Dx=2,Dy=1,ticksize=0 2pt]{->}(0,0)(0,0)(17.5,4.3)
  \psplot[linecolor=black,linewidth=1.2pt]{0}{17}{x sqrt}
\end{pspicture}
\end{center}
\end{minipage}
\begin{minipage}{0.35\linewidth}
Courbe 2:

\begin{center}
\psset{xunit=0.55cm,yunit=0.34cm,arrowsize=2pt 2}
\begin{pspicture}(-3.2,-0.5)(3.2,9.5)
  \psgrid[subgriddiv=4,gridlabels=0pt,gridcolor=gray!35,subgridcolor=gray!18](-3,0)(3,9)
  \psaxes[Dx=1,Dy=2,ticksize=0 2pt]{->}(0,0)(-3.2,0)(3.2,9.4)
  \psplot[linecolor=black,linewidth=1.2pt]{-3}{3}{x 2 exp}
\end{pspicture}
\end{center}

\end{minipage}
\begin{minipage}{0.25\linewidth}
Courbe 3:

\begin{center}
\psset{xunit=0.85cm,yunit=0.48cm,arrowsize=2pt 2}
\begin{pspicture}(-1.4,-4.3)(2.2,7.1)
  \psgrid[subgriddiv=4,gridlabels=0pt,gridcolor=gray!35,subgridcolor=gray!18](-1,-4)(2,7)
  \psaxes[Dx=1,Dy=2,ticksize=0 2pt]{->}(0,0)(-1.3,-4.2)(2.2,7)
  \psplot[linecolor=black,linewidth=1.2pt]{-1.4}{1.92}{x 3 exp}
\end{pspicture}
\end{center}

\end{minipage}

\begin{enumerate}
\item Identifier la fonction de référence représentée par chacune des courbes.
\item En utilisant la courbe appropriée, résoudre graphiquement :

a)~$x^3=3$~\hspace{0.8cm}b)~$x^2<2$~\hspace{0.8cm}c)~$\sqrt{x} \geqslant 2$~\hspace{0.8cm}d)~$x^3>3$~\hspace{0.8cm} e)~$\sqrt{x}<3$~\hspace{0.8cm} f)~$x^2>6$~\hspace{0.8cm}.

\end{enumerate}

\exo \textit{[Chercher]}

\begin{minipage}{13cm}
\begin{enumerate}
\item En utilisant le graphique ci-contre, dire si le nombre admet des antécédents par la fonction carré. Si oui, les déterminer. 

a)~$-1$ \hspace{0,5cm} b)~$4$ \hspace{0,5cm} c)~$0$ \hspace{0,5cm} d)~$-2025$ \hspace{0,5cm} e)~$3$ 

\bigskip

\item En utilisant les résultats précédents, conjecturer le nombre de solutions de l'équation $x^2=a$ en fonction de la valeur de $a$.
\end{enumerate}
\end{minipage}
\hfill
\begin{minipage}{5cm}
\begin{center}
\psset{xunit=0.7cm,yunit=0.7cm,algebraic=true,dimen=middle,dotstyle=o,dotsize=5pt 0,linewidth=1.6pt,arrowsize=3pt 2,arrowinset=0.25}
\begin{pspicture*}(-2.4,-1.5)(2.4,5)
\multips(0,-2)(0,1){7}{\psline[linestyle=dashed,linecap=1,dash=1.5pt 1.5pt,linewidth=0.4pt,linecolor=black]{c-c}(-2.5,0)(2.5,0)}
\multips(-3,0)(1,0){7}{\psline[linestyle=dashed,linecap=1,dash=1.5pt 1.5pt,linewidth=0.4pt,linecolor=black]{c-c}(0,-1.5)(0,5)}
\psaxes[labelFontSize=\scriptstyle,xAxis=true,yAxis=true,Dx=1,Dy=1,ticksize=-2pt 0,subticks=2]{->}(0,0)(-2.4,-1.5)(2.4,5)
\rput{0}(0,0){\psplot[linewidth=2pt]{-5}{5}{x^2/2/0.5}}
\rput[tl](0.5,5){$y=x^2$}
\end{pspicture*}
\end{center}
\end{minipage}



\exo \textit{[Calculer]}

Résoudre les équations suivantes dans $\R$.

a)~$x^2=16$ \hspace{0,5cm} b)~$x^2=-9$ \hspace{0,5cm} c)~$x^2=0$ \hspace{0,5cm} d)~$x^2=\dfrac{36}{100}$ \hspace{0,5cm} e)~$(x+1)^2=0$ \hspace{0,5cm} 


\exo \textit{[Raisonner]}

Margaux affirme : \og Quel que soit le nombre réel $k$, l'équation $\sqrt{x}=k$ admet une solution. \fg{} 

Qu'en pensez-vous?

\exo \textit{[Calculer]}

\begin{enumerate}
\item Déterminer l'antécédent de chaque nombre par la fonction racine carrée.

a)~5~~ \hspace{0.8cm}b)~8~~ \hspace{0.8cm}c)~11~~ \hspace{0.8cm}d)~$\dfrac{3}{2}$~~ \hspace{0.8cm}e)~$10^{-1}$.
\end{enumerate}

\exo \textit{[ Raisonner]}

\begin{enumerate}
\item On souhaite résoudre algébriquement l'inéquation $\sqrt{x} > 1,5$ sur $[0~;~+\infty[$.

\begin{enumerate}[label=\alph*)]
\item Quel est le sens de variation de la fonction carré sur $[0~;~+\infty[$?
\item En utilisant la question précédente, compléter :

$\sqrt{x} > 1,5 \iff ( \vphantom{\sqrt{x}} \ldots)^2 \ldots (\vphantom{1.5} \ldots)^2 \iff x \ldots 2.25 \iff x \in \ldots $

\end{enumerate}
\item Résoudre algébriquement chaque inéquation sur $[0~;~+\infty[$.

 a)~$\sqrt{x} \geqslant \dfrac{1}{4}$~\hspace{0.8cm} b)~$\sqrt{x} \leqslant 2,5$~\hspace{0.8cm} c)~$\sqrt{x} < \dfrac{2}{3}$.
\end{enumerate}



\exo \textit{ [Calculer, Raisonner]}

\bigskip

Voici une fonction, où $a$ est un nombre réel et $a \neq 0$.

\bigskip

\begin{center}
\begin{minipage}{0.46\linewidth}
\begin{pythoncodenumligne}
def Puissance(a):
    n = int(0)
    P = 1
    while P < 100:
        n = n + 1
        P = P*a
    return n
\end{pythoncodenumligne}
\end{minipage}
\end{center}

\medskip

\begin{enumerate}
\item Compléter le tableau de suivi des variables puis donner le résultat renvoyé par \textbf{Puissance(2)}:

\begin{center}
\renewcommand{\arraystretch}{2}
\begin{tabular}{|c|*{9}{p{1.2 cm}|}}
\hline
\centering $Condition P<100$ & \centering X  & \centering Vrai & & & & & & & \tabularnewline
\hline
\centering $n$ & \centering $0$  & \centering $1$ & & & & & & &  \tabularnewline
\hline
\centering $P$ &  \centering $1$ & \centering $2$ & & & & & & &  \tabularnewline
\hline
\end{tabular}
\end{center}

\item  Expliquer ce que renvoie cette fonction.
\end{enumerate}


\exo \textit{ [Raisonner]}

Ranger les nombres suivants par ordre croissant en justifiant la réponse.

a)~$(10^{-2})^3~;~10^{-2}~;~(10^{-2})^2 $~\hspace{0.8cm} b)~$2025^2$~;~$2025^3$~;~$2025$~\hspace{0.8cm} c)~$\dfrac{5}{6}~;~\left(\dfrac{5}{6}\right)^3~;~\left(\dfrac{5}{6}\right)^2 $.


\exo  \textit{ [Raisonner]}

Utiliser la position relative des courbes des fonctions de référence pour comparer les nombres suivants.

\begin{enumerate}
\item $0,604^2 .... 0,604^3$ car la courbe de la fonction ................ est située au dessus de la courbe de la fonction ................ sur l'intervalle ................ .

\bigskip

\item ${\left(\dfrac{7}{8}^3\right)}^2$ .... $\left(\dfrac{7}{8}^3\right)$ car la courbe de la fonction identité est située au ................ de la courbe de la fonction ................ sur l'intervalle ................ .

\bigskip

\item $\pi-3$ .... ${\left(\pi -3\right)}^3$ car la courbe de la fonction identité est située au ................ de la courbe de la fonction ................ sur l'intervalle ................ .

\bigskip

\item ${\left(\sqrt{2}\right)}^2 .... {\left(\sqrt{2}\right)}^3$ car la courbe de la fonction ................ est située au ................ de la courbe de la fonction ................ sur l'intervalle ................ .

\end{enumerate} 

